% tab_main: scripts/cr/analysis/make_tables.py, commit 548eaa68bf12, mode cr, variant default, cell filter completed, levels taskmean
% Generated file, do not edit by hand.
\paragraph{Notes to Table~\ref{tab:main}.} All rows are evaluated on one matched set of $n=440$ (task, seed) cells. It is the intersection of the completed cells of every row except the rows that cover only part of it. A completed cell is a run that made all 25 proposals after the 5 cold starts; runs in which a surrogate fit failed and a step fell back to a uniform proposal are included. The rows that cover only part of the set use the part they cover (GP in nested count-sketch subspaces, Thompson sampling in a trust region: $n=439$, one run crashed; SAAS fully Bayesian GP, qLogEI, L-BFGS-B: $n=104$, 10 seeds). Mean = mean of per-task means; bold marks the highest. $\Delta$ = \geowalk{} minus the row on the same cells, as a mean of per-task means of the paired differences. W/L = cells won/lost by \geowalk{}, ties excluded. $d$ = mean/sd of the cell differences. $p_{\mathrm{BH}}$ (cells) = two-sided Wilcoxon signed-rank test over the cells, Benjamini--Hochberg corrected over the 31 rows. Tasks = tasks, of the 11, on which \geowalk{} has the higher mean. $p_{\mathrm{task}}$ (uncorr.) = exact two-sided Wilcoxon test on the per-task mean differences. It is the primary test of the paper and is printed here without correction. $^\dagger$ Run from the submitted version, not rerun for this version: it ran through a separate harness, or no current environment reproduces it. The $n$ of such a row is the number of its cells inside the matched set. Rows other than the published recipes are named by kernel, acquisition and maximizer, not by the method they adapt. Their deviations from the published methods are summarized in Table~\ref{tab:baselines} and listed in full in the method cards released with the code (\texttt{docs/baselines/}). The trust-region row on the fixed box was added for this version next to the trust-region row of the submitted version. Published recipes follow the authors' released code and carry the authors' names only if they reproduced a published result (Table~\ref{tab:verification}). Their comparison with \geowalk{} and the GP on the shell is Table~\ref{tab:published}. Box rows use the $[p_{01},p_{99}]$ box. Exception: Linear BO (Doumont et al.) runs in the box $[-8,8]^D$ of the authors' released code. A cell with $|\mathrm{diff}|\le10^{-6}$ is a tie: it is excluded from W/L and enters the cell-level Wilcoxon test as a zero difference. Abbreviations: DSP prior is the dimension-scaled lengthscale prior $\mathrm{LogNormal}(\sqrt2+\tfrac12\ln D,\sqrt3)$ of Hvarfner et al. $\ell_0$ is the initial lengthscale before the maximum-likelihood fit. EI, LogEI and qLogEI are expected improvement, its logarithm (analytic) and its Monte Carlo batch version. Candidate-set argmax: the acquisition is maximized over a finite set of random candidates on the support. Sample-around-best: L-BFGS-B restarts include points sampled around the best observations. SAAS is the sparse axis-aligned subspace prior. MSR and MLE (scaled) are the two recipes of Papenmeier et al., with maximum-likelihood lengthscales initialized at $0.1\sqrt D$. MSR adds RAASP perturbations in the acquisition maximizer, and MLE (scaled) does not. The $[p_{01},p_{99}]$ box spans the per-coordinate 1st to 99th percentile of encoded ZINC latents. $\bar r$ is the shell radius, the mean norm of the run's five cold-start latents. $S^{D-1}$ is the unit sphere.
